\chapter{DI and Shit-tuation}

	\noindent
	What does DI have in common with the Shit-tuation context explained in unnecessary comment 006?
	\\\\
	
	\noindent
	First we need a set (an S) with cardinal $\aleph_{k}$. Our S will be a random PACK. Cardinal of our (PACK) S  = $\aleph_{0}$. $K = 0$.
	\\\\
	
	\noindent
	We need injectivity implications:\\\\
	a) If our PACK ends empty after some erradication of lcfs, we don't have nothing.\\\\
	b) If our PACK, after checking ALL possible D-pairs combinations, is NOT EMPTY, it means it keeps some level, so ALL sneis keeps at least one level, so they are all disjoint, and we can apply Naive CA th.
	\\\\
	
	\noindent
	We need the chops, following ALL the conditions of the context.
	\\\\
	
	\noindent
	We have chops: When we quit subpacks from our MEGAPACK, we quit a "concrete list of elements perfectly defined" from the MEGAPACK.
	\\\\
	
	\noindent
	We have $\aleph_{0}$ different possible chops:\\
	SNEIs has cardinal $\aleph_{1}$. (SNEIs X SNEIs) that is the set of ALL possible D-pairs, too. BUT we only have $\aleph_{0}$ different "groups" of D-pairs, classified by their hability to provoke a "chop": Families.
	\\\\
	
	\noindent
	ALL chops are previously defined: System of relations {\ref{tab: system_relations_r_theta_k}; subapcks and levels {\ref{fig:subpacks_levels}}.\\\\
	In that table, you can see ALL the universes you need to quit: all that are equal or smaller than the $\gamma$ value of the D-pair. Defined for ALL possible D-pairs. And you have too, all the Families of D-pairs, you have solved FOREVER, once you quit that universe, and previous ones. When we defined FLJA\_abstract relation, we defined ALL chops in DI too. ALL the infinite, enumerable, list of possible chops.\\\\
	REMEMBER: $F_{0}$ is "solved" when we use the universe $\theta_{1}$. If we have $\gamma=0$ in our D-pair, we don't need to do any erradications. The first universe has index 1, so it is not equal or smaller than 0.
	\\\\
	
	\noindent
	\textbf{Chops are independent}: if you do a chop for a particular D-pair, you don't need the previous chops to make it works.\\\\
	\textbf{All can be executed in parallel}: when you execute one chop, you are executing ALL the previous ones too. WE can talk about what happens when we do all them, because doing one do not interfere with the rest. We are just quitting R-pairs from FLJA\_abstract: if someone did before, what I am trying to do, they just did my job, before me.\\\\
	\textbf{All can be executed in an instant}: We have defined perfectly the list of R-pairs (universes and sub-relations $r_{\theta_k}$). We can define it as an operation of substraction between sets. In mathematics do not exists "time of execution".\\\\
	\textbf{They can ve executed disordered}: Chops in DI are independent, and they do not need an order. That is WHY we talk about a concrete D-pair is a representative of an entire Family. If another D-pair, with the same $\gamma$ level of conflict, made us, previously, to erradicate the universes we need now, nothing happens. The job is done by another dude. No problem. IF we try a D-pair with a smaller $\gamma$... is the same. Other D-pairs can do the job previously, without problem. What we need is the "actual" state of FLJA\_abstract, after the proper erradications, to make our actual D-pair, disjoint (their MEGAPACKs).
	\\\\
	
	\noindent
	\textbf{POINT 3}:\\
	We have in reallity $\aleph_{0}$ different chops, represented by the subpacks, still inside our MEGAPACK, after the erradication provoked by a D-pair that belong to a concrete Family of D-pairs. Let's call each surviving subset of the MEGAPACK: $S_{i}$.
	\\\\
	
	\noindent
	\textbf{Point 3a}:\\
	$S_{i+1} \subset S_{i}$: because they have exactly the same subpacks, minus one.
	\\\\
	
	\noindent
	\textbf{Point 3b}:\\
	$ Card(S_{i} \backslash S_{i+1}) = Card (S)$: The difference is a subpack. A subpack is ALL the lcfs a concrete relation $r_{\theta_{k}}$ assigns to that snei: $\aleph_{0}$. The cardinal of all the MEGAPACK is $\aleph_{0}$: it is a subset of $LCF_{2p}$, which is a subset of LCF, which has a bijection with $\mathbb{N}$.
	\\\\
	
	\noindent
	\textbf{Point 3c}:\\
	For every lcf, you could name, inside our studied MEGAPACK, it belongs to a universe $\theta_{k}$. I can ALWAYS find a D-pair, with a $\gamma$ value bigger than that k, and making you quit that concrete lcf from the MEGAPACK. The infinite intersection of all $S_{i}$ (or all states of the MEGAPACK), is EMPTY.
	\\\\
	
	\noindent
	Now only rest one question: ¿"MEGAPACK" ends empty or not?\\\\
	If it ends NONE EMPTY, well... we have destroyed Cantor's Theorem.\\\\
	If it ends EMPTY, we are saying that under the contenxt of the Shit-tuation, the special set "S" ends TOTALLY EMPTY.
	\\\\
	
	\noindent
	Only one singular lcf surviving inside our MEGAPACK, after ALL erradications of R-pairs ( subpacks in the MEGAPACK ), means we can build the injectivity.\\\\
	It one element survives inside a random MEGAPACK, means the rest of MEGAPACKs have the same condition. We quit them subpack by subpack: if one survives, an entire subpack have survived. And THEN we know they all have cardinal bigger than zero and are ALL DISJOINT between them: we have asked for all combinations of D-pairs. Using the axiom of choice, we can take one lcf, surviving in the MEGAPACK of each possible snei: all different. We have built a proper injectivity.
	\\\\
	
	\noindent
	ALL IS REDUCED TO THIS QUESTION: OUR RANDOM MEGAPACK ("S") ENDS EMPTY OR NOT, under the same context than the Shit-tuation?
	\\\\ 
	
	
	
	
	
	
	
	
	
	
	
	